\documentclass[a4paper,12pt]{article}
\usepackage[T2A]{fontenc}
\usepackage[utf8]{inputenc}
\usepackage[main=russian]{babel}
\usepackage{amsmath,amssymb}
\usepackage{geometry}
\geometry{tmargin=1.5cm,bmargin=1.5cm,lmargin=1.7cm,rmargin=1.7cm}
\setlength{\parindent}{0pt}
\setlength{\parskip}{6pt}

\begin{document}

\begin{center}
{\Large \textbf{Формализация задачи (логика по композитам)}}\\
\vspace{2pt}
{\small (кратко, чтобы переписать на лист)}
\end{center}

\textbf{Дано.}
Выборка объектов \(X=\{x^{(k)}\}_{k=1}^{n}\), где \(x^{(k)}\in\mathbb{R}^{p}\) --- вектор исходных числовых признаков (столбец \texttt{Outcome} при построении не используется).

\textbf{Требуется.}
Получить разбиение на 2 класса \(y^{(k)}\in\{0,1\}\) и построить интерпретируемую логическую модель \(f\), приближающую эти метки по бинарным признакам.

\hrule

\textbf{1) Композитные признаки (сокращение размерности).}
Нужно построить отображение
\[
\Phi:\mathbb{R}^{p}\to\mathbb{R}^{m},\qquad g^{(k)}=\Phi(x^{(k)})=(G_1^{(k)},\dots,G_m^{(k)}),
\]
где \(G_r\) --- агрегат группы зависимых исходных признаков (композит).

\textbf{2) Кластеризация (целевой класс без учителя).}
В пространстве \(g^{(k)}\in\mathbb{R}^{m}\) найти разбиение на 2 кластера, например как решение задачи \(k\)-means:
\[
\min_{\{C_0,C_1\},\,\mu_0,\mu_1}\;\sum_{c\in\{0,1\}}\sum_{k\in C_c}\|g^{(k)}-\mu_c\|^2,
\qquad y^{(k)}=c\ \text{для}\ k\in C_c.
\]

\textbf{3) Разбиение 50/50 и проверка однородности.}
Разделить индексы \(\{1,\dots,n\}=I_{\text{work}}\cup I_{\text{ctrl}}\), \(|I_{\text{work}}|=|I_{\text{ctrl}}|=n/2\).
Требование однородности распределений: для каждого непрерывного признака (используемых \(G_r\)) выполняется критерий Колмогорова--Смирнова при \(\alpha=0.05\):
\[
D_n=\sup_x\bigl|F_{\text{work}}(x)-F_{\text{ctrl}}(x)\bigr|,\qquad
D_{\text{crit}}=1.36\sqrt{\frac{n_1+n_2}{n_1n_2}},\qquad
D_n\le D_{\text{crit}}.
\]

\textbf{4) Бинаризация признаков для логики.}
Для каждого композита задаётся порог \(\theta_r\) и строится бинарный вектор:
\[
b^{(k)}=(b_1^{(k)},\dots,b_m^{(k)})\in\{0,1\}^m,\qquad
b_r^{(k)}=\mathbf{1}\{G_r^{(k)}\le \theta_r\}.
\]

\textbf{5) Логическая модель (СДНФ \(\rightarrow\) минимизированная ДНФ).}
Ищется булева функция \(f:\{0,1\}^m\to\{0,1\}\) в виде ДНФ, согласованная с метками кластеров:
\[
Acc(I)=\frac{1}{|I|}\sum_{k\in I}\mathbf{1}\{f(b^{(k)})=y^{(k)}\},\qquad I\in\{I_{\text{work}},I_{\text{ctrl}}\}.
\]
Построение: (i) СДНФ по положительному классу \(y=1\), (ii) минимизация ДНФ (уменьшение числа импликант) при сохранении согласованности с \(y\).

\textbf{Проверка адекватности (главное) + сравнение с отчётом.}
\begin{itemize}
  \item \textbf{Однородность половин (К--С, \(\alpha=0.05\)).}
  Требование: \(D_n\le D_{\text{crit}}\) для всех используемых непрерывных признаков.
  Для варианта \(m=2\) (композиты \(G1,G2\)) получено:
  \[
  G1:\;D_n=0.0677,\;D_{\text{crit}}=0.0981\;\Rightarrow\;\text{ОК},\qquad
  G2:\;D_n=0.0312,\;D_{\text{crit}}=0.0981\;\Rightarrow\;\text{ОК}.
  \]

  \item \textbf{Согласованность логики с кластеризацией (адекватность интерпретации).}
  На обучающей части: \(Acc(I_{\text{work}})=0.8594\).
  На контрольной части: \(Acc(I_{\text{ctrl}})=0.8724\), матрица ошибок относительно меток \(y\):
  \[
  TP=98,\;FN=0,\;FP=49,\;TN=237.
  \]
  Это означает, что минимизированная ДНФ корректно аппроксимирует разбиение \(k\)-means и может
  рассматриваться как интерпретируемое правило для полученных классов.

  \item \textbf{Сверка с lab3\_report.tex (внешний контроль по Outcome).}
  В отчёте указано, что при \(m=2\) (вариант «по score») внешние метрики ниже, чем при \(m=4\):
  \[
  m=4:\;\mathrm{F1_{macro}}=0.7261,\;\mathrm{BalAcc}=0.7298;\qquad
  m=2:\;\mathrm{F1_{macro}}=0.5842,\;\mathrm{BalAcc}=0.5831.
  \]
  Следовательно, адекватный по смыслу вывод, который должен совпасть с отчётом:
  увеличение числа композитов до \(m=4\) улучшает согласованность с \texttt{Outcome},
  тогда как логическая модель в этой работе проверяется прежде всего на согласованность с метками \(k\)-means.
\end{itemize}

\end{document}

